\TUKEsection{Princíp I²C}{SDA, SCL, pull-up rezistory a open-drain}

\begin{frame}{Ciele prednášky}
\begin{tukebox}{Po tejto prednáške by mal študent vedieť}
\begin{itemize}
    \item vysvetliť základný význam témy: twi/i²c komunikácia,
    \item identifikovať príslušné registre alebo hardvérové bloky,
    \item rozlíšiť abstraktný princíp od konkrétnej implementácie v ATmega328P,
    \item pripraviť jednoduchý program v jazyku C bez Arduino frameworku.
\end{itemize}
\end{tukebox}
\end{frame}

\begin{frame}{Kontext v predmete}
\begin{itemize}
    \item Používame Arduino UNO R3 ako vývojovú dosku.
    \item Cieľovým mikrokontrolérom je ATmega328P.
    \item Programujeme v čistom C s vlastným \codeinline{main()}.
    \item Periférie nastavujeme priamo cez registre.
\end{itemize}
\end{frame}
